\documentclass[11pt,a4paper]{article}
\usepackage[margin=2.2cm]{geometry}
\usepackage{lmodern}
\usepackage[T1]{fontenc}
\usepackage[utf8]{inputenc}
\usepackage{xcolor}
\usepackage{booktabs}
\usepackage{graphicx}
\usepackage{enumitem}
\usepackage{tcolorbox}
\usepackage{titlesec}
\usepackage[colorlinks=true,linkcolor=black,urlcolor=blue]{hyperref}

\definecolor{good}{HTML}{1B7F3B}
\definecolor{warn}{HTML}{B25000}
\definecolor{bad}{HTML}{A31515}
\definecolor{rule}{HTML}{444444}

\titleformat{\section}{\large\bfseries}{\thesection.}{0.6em}{}
\titleformat{\subsection}{\normalsize\bfseries}{\thesubsection}{0.6em}{}
\setlist{nosep,leftmargin=1.4em}

\newtcolorbox{verdict}[1]{colback=#1!6,colframe=#1!70!black,
  boxrule=0.7pt,arc=2pt,left=8pt,right=8pt,top=6pt,bottom=6pt}

% wrapping find/replace rows
\newcommand{\cb}[2]{%
  \par\vspace{2.5pt}\noindent
  \begin{tabular}{@{}p{2.6em}@{}p{0.85\linewidth}@{}}
  {\footnotesize\scshape #1} & {\footnotesize\ttfamily\raggedright #2}
  \end{tabular}\par}
\emergencystretch=3em
\sloppy

\begin{document}

\begin{center}
{\Large\bfseries Integrity Report Review}\\[4pt]
{\normalsize Q\&A-Driven Interactive Storytelling Using Large Language Models:\\
A Multimodal Framework for Personalised Narrative Generation}\\[6pt]
{\small Anjali Madan Jha (124MTCM1008) \quad$\cdot$\quad Guide: Dr.\ Vidyullata Devmane\\
Department of Computer Engineering, Shah \& Anchor Kutchhi Engineering College}\\[6pt]
{\footnotesize Analysed file: \texttt{final paper.pdf} (10 pp., pdfTeX-1.40.29, compiled 27 June 2026)\\
Review date: 18 August 2026}
\end{center}

\vspace{4pt}
\hrule height 0.8pt
\vspace{10pt}

\section*{Summary}
\addcontentsline{toc}{section}{Summary}

Two flags were raised on this paper. They have different answers, and it matters
that they are not treated the same way.

\begin{verdict}{good}
\textbf{Flag 1: Hidden text \quad$\rightarrow$\quad FALSE POSITIVE.}
The 489 white characters are labels printed inside coloured boxes in three
diagrams. Every one of the 60 white text spans sits on a dark filled rectangle
and is plainly legible in the rendered PDF. Nothing is concealed. This flag can
be contested with the evidence in Section 1.
\end{verdict}

\vspace{6pt}

\begin{verdict}{warn}
\textbf{Flag 2: AI writing \quad$\rightarrow$\quad PARTLY REAL.}
Some of this is an artefact of the genre, because methods and results prose is
uniform by nature and detectors score that heavily. But there is a genuine,
repeated stylistic habit in the text, and it is worth fixing. Eight specific
edits are given in Section 3.
\end{verdict}

\vspace{6pt}

\begin{verdict}{bad}
\textbf{Also found, unrelated to either flag: a factual contradiction.}
Section V.J reports ablation numbers that contradict Table V and Section V.F.
This is more damaging than either integrity flag because an examiner comparing
the two will catch it. See Section 4.
\end{verdict}

\section{The hidden-text flag is a false positive}

\subsection*{Where the white text is}

All of it is confined to three figures on three consecutive pages. Pages 1--3
and 7--10 contain no white text at all.

\begin{center}
\begin{tabular}{@{}llrr@{}}
\toprule
\textbf{Page} & \textbf{Figure} & \textbf{Spans} & \textbf{Chars} \\
\midrule
4 & Fig.\ 2, Four-layer system architecture & 28 & 293 \\
5 & Fig.\ 6, Image synthesis pipeline       & 14 & 77 \\
6 & Fig.\ 7, Multimodal audio signal flows  & 18 & 119 \\
\midrule
\multicolumn{2}{@{}l}{\textbf{Total}} & \textbf{60} & \textbf{489} \\
\bottomrule
\end{tabular}
\end{center}

\noindent
The report cites 458 characters; the direct count is 489 across the same three
pages. The small difference is a counting convention, most likely the em dashes
and the \texttt{<img>} angle brackets.

\subsection*{Why it is not concealment}

Each white span was tested against the vector fill shapes beneath it.
\textbf{All 60 of 60 sit on a dark filled box.} Representative fill colours
recovered from the page: RGB(43,\,43,\,43) charcoal, RGB(86,\,0,\,169) purple,
RGB(103,\,26,\,178) violet, RGB(64,\,64,\,191) indigo. White is the only
readable text colour on those fills.

The figures below are rendered directly from the flagged pages at 4$\times$
magnification. The text the report calls hidden is the text you can read here.

\begin{center}
\includegraphics[width=0.92\linewidth]{p4_fig2_arch.png}\\
{\footnotesize Page 4, Fig.\ 2. All 28 flagged spans on this page are these box labels.}
\end{center}

\vspace{4pt}

\begin{center}
\includegraphics[width=0.92\linewidth]{p6_fig7_audio.png}\\
{\footnotesize Page 6, Fig.\ 7. All 18 flagged spans on this page are these box labels.}
\end{center}

\subsection*{Checks for genuine concealment: all negative}

The four techniques actually used to hide text in a PDF were each tested
across all ten pages.

\begin{center}
\begin{tabular}{@{}lc@{}}
\toprule
\textbf{Concealment technique} & \textbf{Instances found} \\
\midrule
Invisible text render mode (\texttt{3 Tr} / \texttt{7 Tr} operators) & 0 \\
Sub-legible font size ($<$ 4\,pt)                                    & 0 \\
Text positioned outside the visible page rectangle                   & 0 \\
Text covered by an opaque image                                      & 0 \\
\bottomrule
\end{tabular}
\end{center}

\subsection*{Complete colour census}

Every text character in the document, grouped by colour. There is no fourth
category and no white body paragraph.

\begin{center}
\begin{tabular}{@{}lrl@{}}
\toprule
\textbf{Colour} & \textbf{Chars} & \textbf{What it is} \\
\midrule
RGB(0,\,0,\,0)       & 42{,}669 & Body text, captions, tables \\
RGB(255,\,255,\,255) & 489      & Diagram box labels (Figs.\ 2, 6, 7) \\
RGB(133,\,133,\,133) & 26       & Figure annotation (``Modular Core'') \\
RGB(115,\,115,\,115) & 8        & Figure annotation \\
\bottomrule
\end{tabular}
\end{center}

\noindent
There is no keyword stuffing, no inflated word count, and no text that is
invisible to a reader but visible to a similarity checker.

\subsection*{Recommended response}

Contest the flag. One sentence is enough: \emph{the white characters are labels
inside coloured boxes in Figures 2, 6 and 7, and they are visible in the
rendered PDF.} Point the reviewer at those three figures. Nothing in the paper
needs to change.

If the committee requires the flag cleared regardless, compile each of the three
diagrams as a standalone PDF and convert its text to vector outlines
(\texttt{gs -o fig-flat.pdf -dNoOutputFonts -sDEVICE=pdfwrite fig.pdf}), then
include the flattened file. This removes the text layer. It is worth noting this
makes the figures slightly worse for screen readers in order to satisfy a
detector that has misread them.

\section{The AI-writing flag}

\subsection*{What is genre artefact}

The flagged region is Sections IV.E through V.F: endpoint names, table
pointers, percentages. This prose is uniform and low-entropy by nature, which
is what AI detectors measure. Constructions such as ``Three evaluation axes were
used'' or ``Table IV reports mean latencies'' are correct IEEE register. Passive
voice and formulaic figure pointers are how methods sections are written, and no
rewrite makes them read as less mechanical to a classifier. These tools have
well-documented false-positive rates on exactly this genre.

\subsection*{What is genuinely worth fixing}

The following patterns are real and repeat throughout the paper, not only in the
flagged region.

\begin{center}
\begin{tabular}{@{}lcl@{}}
\toprule
\textbf{Pattern} & \textbf{Count} & \textbf{Assessment} \\
\midrule
Participial clause-tails (\texttt{, confirming\ldots}) & 13 & Dominant tell \\
Em dashes                                              & 20 & Secondary \\
Evaluative filler in results sections                  & 4  & Editorialising \\
Caption text duplicating body text                     & Figs.\ 9, 10 & Also a reviewer complaint \\
\bottomrule
\end{tabular}
\end{center}

\noindent
The strongest signal is clustering: Sections V.D, V.E and V.F each close on a
participial tail (\emph{confirming\ldots}, \emph{indicating\ldots},
\emph{confirming\ldots}). Three in a row across consecutive subsections is what
reads as machine-generated. Removing them is the single highest-value change.

\section{Eight edits}

Facts and figures are unchanged. Register stays formal, because plain and
neutral is the correct human voice for an IEEE paper. These apply directly in
the Overleaf editor.

\begin{enumerate}[label=\textbf{\arabic*.}]
\item \textbf{Section V.D}
\cb{find}{, confirming Q\&A elicitation as the core value driver}
\cb{repl}{. Both derive directly from the Q\&A step}

\item \textbf{Section V.E}
\cb{find}{, indicating a need for improved discoverability cues}
\cb{repl}{. Both sit behind a secondary tab, which is the most likely explanation}

\item \textbf{Section V.F}
\cb{find}{, confirming that deeper Q\&A elicitation adds measurable value}
\cb{repl}{(delete it; the numbers already make the point)}

\item \textbf{Section IV.E}
\cb{find}{(enabling session replay)}
\cb{repl}{, which is what makes session replay possible}
\cb{find}{(persisting provider and temperature settings)}
\cb{repl}{, which holds the provider and temperature the user selected}

\item \textbf{Section V.D}
\cb{find}{disaggregates satisfaction}
\cb{repl}{breaks satisfaction down}

\item \textbf{Section V.B} (removes the em-dash pair). In the \texttt{.tex}
      the em dash is typed as three hyphens, shown below as {\ttfamily -}{\ttfamily -}{\ttfamily -}.
\cb{find}{1.2 s continuation{-}{-}{-}within the 5 s design target{-}{-}{-}with continuation faster}
\cb{repl}{1.2 s continuation, both inside the 5 s design target. Continuations are faster}

\item \textbf{Section V.J \quad$\cdot$\quad highest priority} (see Section 4)
\cb{find}{improves personalisation by $\approx$ 0.95 Likert points while coherence stays flat}
\cb{repl}{improves personalisation by $\approx$1.4 Likert points, with coherence rising more modestly (0.7 points)}

\item \textbf{Captions of Figs.\ 9 and 10.} Rewrite so they describe what the
      axes show rather than restating their paragraphs. At present the same
      sentence appears twice on one page. For example:
\cb{repl}{Fig. 10. Per-feature adoption across 15 user-study sessions. Each bar is the share of sessions in which the feature was used at least once.}
\end{enumerate}

\section{Priority: Section V.J contradicts Table V}

This is independent of both integrity flags and should be fixed first.

Table V (question-budget ablation) gives these values:

\begin{center}
\begin{tabular}{@{}lcccc@{}}
\toprule
\textbf{Budget} & \textbf{Coherence} & \textbf{Creativity} & \textbf{Personalisation} & \textbf{Time} \\
\midrule
4 (Quick)     & 3.6 & 3.8 & 3.2 & 45\,s \\
8 (Balanced)  & 3.9 & 4.1 & 3.8 & 1.5\,m \\
12 (Detailed) & 4.2 & 4.2 & 4.3 & 2.5\,m \\
17 (Full)     & 4.3 & 4.1 & 4.6 & 4.6\,m \\
\bottomrule
\end{tabular}
\end{center}

\noindent
\textbf{Section V.F reports this correctly:} personalisation rises
$\approx$1.4 points (3.2 $\rightarrow$ 4.6) and coherence rises
3.6 $\rightarrow$ 4.3.

\vspace{4pt}
\noindent
\textbf{Section V.J then states:} ``improves personalisation by $\approx$ 0.95
Likert points while coherence stays flat.''

\vspace{4pt}
\noindent
Both halves are wrong against the paper's own table. The personalisation gain is
1.4, not 0.95. Coherence rises 0.7 points, which is not flat. Correct V.J as
given in edit 7.

\vspace{6pt}
\noindent
\textbf{Secondary observation.} Creativity falls from 4.2 to 4.1 between the
12-question and 17-question budgets. The text does not mention this, but it
supports the existing recommendation of 12 questions as the default and is worth
a clause.

\section{Status of the source files}

The \texttt{.tex} source for the reviewed paper is not on this machine. It
exists here only as a compiled PDF. Two older copies of the paper source were
located:

\begin{center}
\begin{tabular}{@{}lccc@{}}
\toprule
\textbf{Check} & \textbf{projects pack copy} & \textbf{Anjali maam/latex} & \textbf{Flagged PDF} \\
\midrule
Lines                      & 898 & 909 & --- \\
Title                      & ``Q\&A-\emph{Based}\ldots'' & ``Q\&A-\emph{Based}\ldots'' & ``Q\&A-\emph{Driven}\ldots'' \\
\texttt{tikzpicture}       & 0 & 0 & 3 \\
\texttt{core value driver} & 0 & 0 & present \\
\texttt{0.95 Likert}       & 0 & 0 & present \\
CoNoder / Salmaze refs     & 0 & 0 & cited \\
\bottomrule
\end{tabular}
\end{center}

\noindent
Searches covered every \texttt{.tex} and \texttt{.bib} in the home directory,
plus the interior of every \texttt{.docx}, \texttt{.pptx} and \texttt{.zip} on
the Desktop. None contain the flagged sentences or references [32]--[42].

The reviewed version therefore adds a rewritten related-work section, three
TikZ diagrams and a changed title relative to both local copies. To apply these
edits to the real paper, export the source from Overleaf
(Menu $\rightarrow$ Download $\rightarrow$ Source).

\vspace{6pt}
\noindent
\textbf{Caution.} Do not rebuild from the local copies. Doing so loses the newer
related-work section and the TikZ figures, and would submit a paper whose title
differs from the one that was reviewed.

\section{A note on disclosure}

Whether any part of this text was drafted with AI assistance cannot be
determined from the file, and is not assumed here. If it was, editing the
stylistic markers does not substitute for whatever disclosure the institution
requires, and following that policy is the safer course than reducing a detector
score. If it was not, the eight edits above are ordinary copy-editing that
improve the paper, and the caption duplication and the Table V contradiction
were worth catching either way.

\vfill
\hrule height 0.4pt
\vspace{4pt}
{\footnotesize Evidence for Section 1 was produced with PyMuPDF 1.26.5 over
\texttt{final paper.pdf}: per-span colour extraction, vector fill-shape
intersection at each span centre, content-stream scan for text render-mode
operators, and 4$\times$ region rendering. Figure crops are unmodified renders
of pages 4 and 6.}

\end{document}
